\documentclass[../dayan_kernel_manual.tex]{subfiles}

\begin{document}

\section{标准库}\label{sec:kernel-040}

Dayan 提供一组 freestanding C 头文件。它们使用标准名称，但只实现本章列出的接口；
没有文件系统、输入流、时间、locale 或完整 libm。

\subsection{输出与格式化}

\begin{dayantable}{\code{stdio.h}}{tab:kernel-040}{X[24,l] X[76,l]}
  符号 & 行为 \\
  \code{putchar} & 向当前 proc 输出缓冲追加一个字节；失败返回 \code{EOF}。 \\
  \code{puts} & 追加字符串和换行；失败返回 \code{EOF}。 \\
  \code{fflush} & 刷新 \code{stdout} 或空 stream；其它 stream 返回 \code{EOF}。 \\
  \code{printf} & 写入 proc 输出缓冲，缓冲满时通过 RPC 发送。 \\
  \code{vprintf} & 接收 \code{va_list} 的 printf。 \\
  \code{snprintf} & 写本地内存，遵循 size 限制并返回完整输出所需长度。 \\
  \code{vsnprintf} & 接收 \code{va_list} 的 snprintf。 \\
\end{dayantable}

格式化支持 \code{\%\%}、\code{\%c}、\code{\%s}、有符号/无符号十进制、十六进制
和指针；支持 \code{-}/\code{0} flag、十进制宽度以及 \code{l}/\code{ll}/\code{z}
length。不支持浮点、precision、八进制或 \code{\%n}。

输出缓冲由 proc 创建者提供。容量为 0 时，流式输出失败；\code{snprintf} 不受此限制。
proc 退出会尝试刷新，但无法向 kernel 返回失败。需要确认输出送达时，在返回前检查
\code{fflush(stdout)}。

\subsection{私有堆}

\begin{dayantable}{\code{stdlib.h}}{tab:kernel-041}{X[24,l] X[76,l]}
  符号 & 行为 \\
  \code{malloc} & 从当前 hart 私有堆分配至少 16-byte 对齐的块；size 为 0 或失败时返回空。 \\
  \code{calloc} & 检查乘法溢出，并把分配范围清零。 \\
  \code{realloc} & 支持空 ptr 和 size 为 0；扩展时分配、复制、释放，不原地扩大。 \\
  \code{free} & 释放当前 hart 的私有堆块；空指针无操作。 \\
  \code{abort} & 停止当前 hart 的 Dayan OS，不返回。 \\
  \code{exit} & 忽略 status，停止当前 hart 的 Dayan OS，不返回。 \\
\end{dayantable}

地址必须来自同一 hart 的 allocator。不能跨 hart 释放，也不能把 coherent allocation
交给 \code{free}/\code{realloc}。kernel 正常返回只结束当前 proc；通常不应调用
\code{exit} 或 \code{abort}。

\subsection{内存与字符串}

\begin{dayantable}{\code{string.h}}{tab:kernel-042}{X[25,l] X[75,l]}
  符号 & 行为 \\
  \code{memcpy} & 复制不重叠的字节区间。 \\
  \code{memmove} & 复制可重叠的字节区间。 \\
  \code{memset} & 用 value 的低 8 位填充。 \\
  \code{memcmp} & 按 unsigned byte 比较。 \\
  \code{strlen}、\code{strnlen} & 计算字符串长度；后者限制最大检查长度。 \\
  \code{strcpy}、\code{strncpy} & 复制字符串；源过长时 strncpy 可能不写终止 NUL。 \\
  \code{strcmp}、\code{strncmp} & 比较 NUL 结尾字符串。 \\
  \code{strchr}、\code{strrchr} & 查找字符第一次或最后一次出现的位置。 \\
\end{dayantable}

这些函数是标量实现，不提供 DMA 优化或跨核内存顺序。

\subsection{单精度数学}

\begin{dayantable}{\code{math.h}}{tab:kernel-043}{X[24,l] X[76,l]}
  符号 & 行为 \\
  \code{fabsf} & 清除符号位。 \\
  \code{copysignf} & 复制符号位。 \\
  \code{fminf}、\code{fmaxf} & 按 C 语义处理单个 NaN 和有符号零。 \\
  \code{sqrtf} & 使用 RISC-V 单精度平方根指令。 \\
  \code{rsqrtf} & 位级初值加三轮 Newton 迭代的倒平方根近似。 \\
  \code{expf} & 单精度 range reduction 与多项式近似。 \\
  \code{logf} & 单精度 range reduction 与级数近似。 \\
\end{dayantable}

头文件还定义 \code{INFINITY}、\code{NAN}、\code{HUGE_VALF} 和浮点分类宏。没有
double 版本；\code{rsqrtf}/\code{expf}/\code{logf} 不保证通用正确舍入。

\subsection{基础头文件}

\code{stdint.h}、\code{stddef.h} 和 \code{limits.h} 提供当前 ABI 所需的类型与边界；
\code{stdarg.h} 和 \code{stdbool.h} 封装编译器的可变参数与布尔类型。它们不是完整的
Hosted C 标准头，例如没有 \code{PRI*} 宏族。

\end{document}
